\documentclass[11pt, a4paper]{article}

\usepackage[utf8]{inputenc}
\usepackage{amsmath, amssymb, amsthm}
\usepackage{geometry}
\usepackage{hyperref}
\usepackage{xcolor}
\usepackage{braket}
\usepackage{graphicx}

\geometry{margin=1in}

\newtheorem{theorem}{Theorem}
\newtheorem{proposition}{Proposition}
\newtheorem{definition}{Definition}
\newtheorem{remark}{Remark}

\newcommand{\Hgen}{\hat{\mathcal{H}}}
\newcommand{\Kry}{\mathcal{K}}
\newcommand{\Rset}{\mathbb{R}}
\newcommand{\Cset}{\mathbb{C}}
\newcommand{\Eblr}{\mathbb{E}}

\title{\textbf{Analytical Quantum Flow Matching via Inversion-Free
Krylov Reduction in the \texttt{unfer} Kernel}\\[4pt]
\large A Self-Contained, Pedagogical Account}
\author{Implementation Architecture Report}
\date{June 2026}

\begin{document}

\maketitle

\begin{abstract}
We give a self-contained account of a scalable, \emph{analytical} Quantum
Flow Matching (QFM) module for the \texttt{unfer} probability kernel. Recent
work on QFM (Cui, Zhang, and Tang, 2025) shows that mapping generative flows
onto continuous quantum evolution is a powerful alternative to classical
diffusion. That construction, however, relies on Parameterized Quantum
Circuits (PQCs) trained by heuristic variational algorithms, which suffer
from barren plateaus and heavy sampling overhead. We replace the trained
circuit with a closed-form generator. The two ingredients are classical: (i)
data points are encoded as mutually orthogonal single-particle states of a
symmetric Fock space, so that the flow potential \emph{decouples} and the
training reduces to a diagonal system solvable in $\mathcal{O}(M)$ for $M$
data points; and (ii) the resulting bounded generator is compressed by the
Shift-Invert Rational Krylov (SIRK) method of Hashimoto and Nodera, which
gives exponential approximation-error decay and reduces single-step inference
to an $\mathcal{O}(m^2)$ matrix exponential of dimension $m\ll M$. We also
make explicit the measure-theoretic foundation of the construction: in the
infinite-dimensional limit the uniform measure on a hypersphere becomes the
Gaussian measure, and the Gegenbauer polynomials that diagonalize it become
the Hermite polynomials (López and Temme, 1999; Mehler, 1866). This
identifies the ``no-information'' prior with the Fock vacuum and any orthogonal
excitation with acquired information. We close by detailing the integration
into \texttt{unfer}'s GPU-accelerated \texttt{fock\textunderscore sirk}
solver, its \texttt{nested\textunderscore fock\textunderscore algebra} CAS,
and the Austral module runtime.
\end{abstract}

\tableofcontents

\section{Introduction}

\subsection{Flow matching, classical and quantum}

\emph{Flow matching} has become a dominant paradigm in classical generative
modeling. The idea is geometric: rather than learning to denoise, one learns
a deterministic velocity field $u_t(x)$ whose ordinary differential equation
(ODE) $\dot{x}=u_t(x)$ transports a simple prior distribution $p_0$ (typically
Gaussian noise) into a complex target distribution $p_1$ (the data) along a
prescribed interpolation. The training objective is a regression: match the
network's velocity to a known \emph{conditional} target velocity, whose
expectation is the true marginal velocity. Because the target is available in
closed form along each interpolation path, training avoids simulating the ODE
and is highly efficient.

Cui, Zhang, and Tang (2025) extend this paradigm to the quantum domain as
\emph{Quantum Flow Matching} (QFM)~\cite{Cui2025}. There the objects
transported are not probability densities but density matrices $\rho$, and the
flow is realized by a quantum circuit that learns a density-matrix propagator
interpolating between an initial ensemble $S_0=\{\ket{\psi_0^m}\}_{m=1}^M$ and
a target ensemble $S_{\mathcal{T}}$. The circuit alternates unitary layers
with measured ancilla qubits, so that a \emph{single} fixed circuit can emulate
the entire density-matrix dynamics --- a decisive advantage on quantum
hardware, where redesigning circuits per state is the dominant cost. QFM
contains the quantum denoising diffusion model (QuDDPM) as the special case in
which the initial ensemble is Haar-random.

\subsection{The variational bottleneck}

The practical QFM circuit of~\cite{Cui2025} is an Entanglement-varied
Hardware-efficient Ansatz (EHA): each layer applies single-qubit rotations
$U^{(1)}=\prod_{j\in\{x,y,z\}}\exp(-i\theta_j\sigma_j/2)$ and entangling gates
$U^{(2)}=\prod_{j}\exp(-i\theta_j\,\sigma_j\otimes\sigma_j/2)$, with the
parameters $\vec{\theta}_\tau$ fit by gradient descent against a fidelity loss.
This inherits the well-known pathologies of variational quantum algorithms:
susceptibility to \emph{barren plateaus} (exponentially vanishing gradients),
convergence to poor local minima, and a sampling-heavy classical--quantum
optimization loop whose cost grows with the ensemble size $M$. The authors
themselves bound the generalization error of a trained layer with
$\mathcal{N}$ gates by
\begin{equation}
  R_P(\vec\theta^{\,\mathrm{op}}_\tau)\;\le\;
  2D_\tau(\vec\theta^{\,\mathrm{op}}_\tau)
  +\mathcal{O}\!\left(\sqrt{\tfrac{\mathcal{N}\log\mathcal{N}}{M}}\right),
  \label{eq:genbound}
\end{equation}
which is reassuring asymptotically but says nothing about reaching a small
training loss $D_\tau$ in the first place --- precisely where the plateaus bite.

\subsection{This work: an analytical alternative}

We propose a deliberate departure from heuristic PQC optimization: an
\emph{analytical} QFM whose generator is written in closed form, with no
neural network and no variational loop. The construction rests on two
classical pillars, developed pedagogically in
Sections~\ref{sec:flowprimer}--\ref{sec:krylov}:

\begin{enumerate}
\item \textbf{Orthogonal Fock encoding (Section~\ref{sec:framework}).} Encoding
each data point as a distinct single-boson mode makes the data potential
\emph{diagonal}. The flow-matching normal equations then decouple into $M$
independent scalar equations, so ``training'' is an $\mathcal{O}(M)$
closed-form average with zero data loss.

\item \textbf{Inversion-free Krylov reduction
(Section~\ref{sec:krylov}).} The resulting time-independent generator
$\bar H$ is compressed by the Shift-Invert Rational Krylov (SIRK) method of
Hashimoto and Nodera~\cite{Hashimoto2019}, which enjoys exponential
error decay in the subspace dimension $m$. Because the potential has no
cross-terms, each Krylov step costs $\mathcal{O}(M)$, and inference is a
single $\mathcal{O}(m^2)$ matrix exponential.
\end{enumerate}

Section~\ref{sec:mehler} supplies the measure-theoretic bridge that makes the
``Gaussian prior $=$ Fock vacuum'' identification precise, via the
Mehler/López--Temme limit of Gegenbauer polynomials to Hermite polynomials.
Section~\ref{sec:impl} maps every piece onto the existing \texttt{unfer}
modules. Throughout we keep prerequisites minimal: a reader comfortable with
ladder operators, Gaussian integrals, and the matrix exponential should be
able to follow without external references.

\section{Background I: Operator functions and the SIRK method}
\label{sec:flowprimer}

QFM ultimately asks us to apply $e^{-iHt}$ (or, more generally, an operator
$\varphi$-function) to a state. We summarize the numerical machinery that
makes this tractable for the large, ``unbounded-like'' generators that arise
from field theories, following Hashimoto and Nodera~\cite{Hashimoto2019}.

\subsection{Exponential integrators and \texorpdfstring{$\varphi$}{phi}-functions}

Evolution equations $u'=Au+g$ have the variation-of-constants solution
\begin{equation}
  u(t)=e^{tA}u_0+\int_0^t e^{(t-s)A}g(s)\,\mathrm{d}s .
\end{equation}
\emph{Exponential integrators} discretize the integral using the family of
$\varphi$-functions
\begin{equation}
  \varphi_0(z)=e^z,\qquad
  \varphi_k(z)=\int_0^1 e^{sz}\,\frac{(1-s)^{k-1}}{(k-1)!}\,\mathrm{d}s
  \quad(k\ge 1),
\end{equation}
e.g.\ the one-step scheme $u_i=e^{\delta A}u_{i-1}+\delta\,\varphi_1(\delta
A)\,g_{i-1}$. The computational kernel is therefore the product
$\varphi_k(A)\,v$ of an operator function with a vector --- exactly the
primitive QFM inference needs (with $k=0$, $A=-iH$).

\subsection{Why naive (Arnoldi) Krylov fails for stiff generators}

The Arnoldi method builds the polynomial Krylov subspace
\begin{equation}
  \Kry_m(A,v)=\mathrm{span}\{v,Av,\dots,A^{m-1}v\},
\end{equation}
and approximates $\varphi_k(A)v\approx V_m\,\varphi_k(\mathbf{H}_m)\,V_m^*v$,
where $V_m$ holds an orthonormal basis and
$\mathbf{H}_m=V_m^*AV_m$ is the compression of $A$. For a \emph{bounded} $A$
the error decays geometrically, governed by polynomial approximation of
$\varphi_k$ on the numerical range $W(A)$. But generators from differential
operators are effectively \emph{unbounded}: $W(A)$ is unbounded, the relevant
conformal map ceases to exist, and the error is no longer guaranteed to decay
in $m$. This is the stiffness obstruction.

\subsection{Resolvents, rational Krylov, and SIRK}

The cure is to work with the \emph{resolvent} $X=(\gamma I-A)^{-1}$, which is
bounded whenever $\gamma$ is outside the spectrum of $A$. One can define
operator functions of an unbounded $A$ through the bounded $X$ via
$f(A):=f_\gamma(X)$ with $f_\gamma(z)=f(\gamma-z^{-1})$. The
\emph{Rational Krylov} subspace uses a sequence of shifts $\gamma_j$ and
resolvents $X_j=(\gamma_j I-A)^{-1}$,
\begin{equation}
  \mathcal{Q}_m(\{X_j\},v)
  =\mathrm{span}\{v,\,X_1v,\,X_2X_1v,\dots,X_{m-1}\cdots X_1 v\}.
  \label{eq:ratkrylov}
\end{equation}
The Shift-Invert Rational Krylov (SIRK) method chooses the real, equispaced
shifts $\gamma_j=N-hj$ with $N>hj>0$. A key algebraic identity is that
all the resolvents share a single ``base'' resolvent $X_m$,
\begin{equation}
  X_j=(\gamma_j I-A)^{-1}=\bigl(I+h(m-j)X_m\bigr)^{-1}X_m,
\end{equation}
so that the rational Krylov space~\eqref{eq:ratkrylov} equals
$\{r(X_m)v:\,r\in\mathcal{R}^{\mathrm{SIRK}}_{m-1}\}$ with denominators
$q_m(z)=(1+hmz)\cdots(1+hz)$. Hashimoto and Nodera prove an error bound that,
crucially, decays \emph{exponentially}:
\begin{equation}
  \bigl\|\varphi_k(A)v-V_m\,\psi_{k,\gamma_m}(\mathbf{H}_m\mathbf{K}_m^{-1})\,V_m^*v\bigr\|
  \;\le\;2C\,\|v\|\,e^{-hm}\!\!\min_{r\in\mathcal{R}^{\mathrm{SIRK}}_{m-1}}\!
  \|f_{k,N}-r\|_{\infty,\Sigma},
  \label{eq:sirkbound}
\end{equation}
with $C\in[2,11.08]$ and $\Sigma$ a bounded convex hull of resolvent numerical
ranges. The factor $e^{-hm}$ is the practical payoff: where Arnoldi (for
unbounded $A$) and even plain rational Krylov decay only polynomially, SIRK
decays geometrically, so a small subspace dimension $m$ (e.g.\ $m=15$--$20$)
suffices. This is the engine \texttt{unfer} already ships as
\texttt{fock\textunderscore sirk}; Section~\ref{sec:krylov} explains why our
decoupled QFM generator makes it run especially fast.

\section{Background II: from flow matching to a continuity Hamiltonian}

We now set up the analytical flow. A flow-matching model transports a prior
$p_0$ to data $p_1$ along the straight-line interpolation
$x_t=(1-t)x_0+t\,x^{(k)}$, whose conditional optimal-transport velocity is the
constant $u_t(x\mid x_0,x^{(k)})=x^{(k)}-x_0$. We will represent the evolving
distribution by a wavefunction $\Psi_t(x)$ on configuration space and derive a
\emph{generator} whose evolution reproduces the flow. The remainder of the
paper makes each object concrete; the rest of this section fixes notation.

A conservative velocity field is the gradient of a scalar potential,
$u_t(x)=-\nabla V_t(x)$. The continuity equation
$\partial_t p_t+\nabla\!\cdot(p_t u_t)=0$ for the probability current then
translates, at the level of the wavefunction amplitude, into a first-order
\emph{continuity Hamiltonian}
\begin{equation}
  \Hgen_t=H_0-\tfrac{i}{2}\bigl(\nabla V_t(x)\cdot\nabla+\nabla^2 V_t(x)\bigr),
  \label{eq:continuity}
\end{equation}
where $H_0$ encodes the prior (Section~\ref{sec:framework}). The two
derivative terms are the drift ($\nabla V\cdot\nabla$) and the
compressibility/source ($\nabla^2 V$) of the flow. Everything hinges on
choosing the potential $V_t$ so that (a) the optimization for its parameters
is analytic and (b) the operators commute in time. The orthogonal Fock
encoding achieves both.

\section{Theoretical Framework: orthogonal Fock encoding}
\label{sec:framework}

\subsection{Zero-loss data encoding and the symmetric Fock space}

To avoid the $\mathcal{O}(M^3)$ matrix inversion that plagues analytical kernel
methods, we treat the $M$ training points as events of a Poisson counting
measure. Via the Wiener--It\^o--Segal isomorphism, each data point
$x_j\in\Rset^D$ is mapped to a \emph{strictly orthogonal} single-particle
state of a symmetric Fock space:
\begin{equation}
  \ket{x_j}=\hat c_j^\dagger\ket{0},\qquad
  \braket{x_i | x_j}=\delta_{ij}.
\end{equation}
In the coordinate representation these correspond to localized, normalized
wave-packets $\Psi_j(x)=\braket{x | x_j}$ centered at $x_j$. Because the
events of a point process are disjoint, the packets have \emph{disjoint
supports}, which gives the two identities the whole construction exploits:
\begin{align}
  \Psi_j(x)\,\Psi_k(x)&=\delta_{jk}\,\Psi_j(x)^2,
  \label{eq:disjoint}\\
  \nabla\Psi_j(x)\cdot\nabla\Psi_k(x)&=0\quad(j\neq k).
  \label{eq:graddisjoint}
\end{align}
The orthogonality~\eqref{eq:disjoint} is ``zero data loss'': no two data points
share amplitude. The gradient identity~\eqref{eq:graddisjoint} is what
decouples the optimization.

\subsection{Analytical potential optimization}

Write the flow potential as a linear combination of the orthogonal packets,
\begin{equation}
  V_t(x)=\sum_{j=1}^M\alpha_j(t)\,\Psi_j(x).
\end{equation}
Fitting the conservative velocity $-\nabla V_t$ to the target
$u_t=x^{(k)}-x_0$ is the least-squares problem
\begin{equation}
  \mathcal{L}(\vec\alpha(t))=
  \Eblr_{t,x_0,x^{(j)},x_t}\Bigl\|
  -\sum_{j=1}^M\alpha_j(t)\,\nabla\Psi_j(x_t)-(x^{(j)}-x_0)\Bigr\|_2^2 .
\end{equation}
Setting $\partial\mathcal{L}/\partial\alpha_k=0$ gives the normal equations
\begin{equation}
  \sum_{j=1}^M\alpha_j(t)\,
  \Eblr_{x_t}\!\bigl[\nabla\Psi_j(x_t)\cdot\nabla\Psi_k(x_t)\bigr]
  =-\,\Eblr_{x_t,x_0,x^{(j)}}\!\bigl[(x^{(j)}-x_0)\cdot\nabla\Psi_k(x_t)\bigr].
\end{equation}
By the gradient-disjointness~\eqref{eq:graddisjoint} the left-hand matrix is
\emph{strictly diagonal}, and along any path $j\neq k$ the packet
$\Psi_k(x_t^{(j)})$ vanishes. The system collapses to a closed-form, per-mode
solution:
\begin{equation}
  \alpha_k(t)=-\,\frac{\Eblr_{x_0}\!\bigl[(x^{(k)}-x_0)\cdot
  \nabla\Psi_k(x_t^{(k)})\bigr]}
  {M\,\Eblr_{x_0}\!\bigl[\|\nabla\Psi_k(x_t^{(k)})\|^2\bigr]},
  \label{eq:alpha}
\end{equation}
with $x_0\sim\mathcal{N}(0,\mathbf{I})$ the Gaussian (Mehler) noise prior and
$x_t^{(k)}=(1-t)x_0+t\,x^{(k)}$. Each $\alpha_k$ is an independent scalar
average, so the whole ``training'' is $\mathcal{O}(M)$ with zero data loss ---
no plateaus, no inner optimization loop.

\subsection{The continuity Hamiltonian and the Mehler prior}

We take the prior to be the rank-1 projector onto the Mehler ground state,
\begin{equation}
  H_0=\ket{0}\bra{0},
\end{equation}
the maximum-entropy state representing total prior uncertainty (made precise in
Section~\ref{sec:mehler}). Substituting $V_t=\sum_j\alpha_j(t)\Psi_j$ into the
continuity Hamiltonian~\eqref{eq:continuity} yields
\begin{equation}
  \Hgen_t=\ket{0}\bra{0}-\frac{i}{2}\sum_{j=1}^M\alpha_j(t)\,
  \hat h_j(x,\nabla),\qquad
  \hat h_j=\nabla\Psi_j\cdot\nabla+\nabla^2\Psi_j .
  \label{eq:Hgen_full}
\end{equation}
For any out-of-distribution state $\ket{x_{\mathrm{OOD}}}$, $\Hgen_t$ first
projects onto the Mehler prior $\ket{0}$, which then distributes amplitude
\emph{uniformly} to the data channels $\ket{x_j}$ --- generalization is built
in, and the construction never touches a spatial grid.

\subsection{Time-averaging and exact commutativity}

Because the packets have disjoint supports, the localized operators act on
independent regions and \emph{exactly commute}:
\begin{equation}
  [\hat h_j,\hat h_k]=0\quad\forall j,k
  \;\Longrightarrow\;[\Hgen_t,\Hgen_{t'}]=0 .
\end{equation}
The time-ordering operator therefore drops out of the evolution, and we may
analytically time-average the parameters, $\bar\alpha_j=\int_0^1\alpha_j(t)\,
\mathrm{d}t$. The continuous-time flow collapses to a single,
\emph{time-independent} generator,
\begin{equation}
  \bar H=\ket{0}\bra{0}-\frac{i}{2}\sum_{j=1}^M\bar\alpha_j\,\hat h_j ,
  \label{eq:Hbar}
\end{equation}
and inference is one exact step, $\Psi_{t=1}=e^{-i\bar H}\Psi_0$.

\paragraph{The implemented (diagonal) generator.}
The differential operators $\hat h_j$ of~\eqref{eq:Hbar} mix the vacuum with
the data channels; that off-diagonal generator is the full QFM flow. The
\texttt{unfer} builtin (Section~\ref{sec:impl}) currently implements the
\emph{diagonal} number-operator surrogate
\begin{equation}
  \bar H_{\text{impl}}=\ket{0}\bra{0}+\sum_{j=1}^M\bar\alpha_j\,\hat n_j,
  \qquad \hat n_j=\hat c_j^\dagger\hat c_j ,
  \label{eq:Hdiag}
\end{equation}
for which the vacuum and each data channel are eigenstates
($\bar H_{\text{impl}}\ket{0}=\ket{0}$,
$\bar H_{\text{impl}}\ket{x_j}=\bar\alpha_j\ket{x_j}$). This is the
spectrally-stable skeleton of~\eqref{eq:Hbar}: it shares the prior, the
$\mathcal{O}(M)$ construction, and the SIRK compressibility, while the
off-diagonal mixing $\hat h_j$ is a documented future extension.

\section{Relating the Fock space with a uniform measure in an
infinite-dimensional sphere}
\label{sec:mehler}

This section justifies, from first principles, the central modeling choice:
\emph{the Gaussian (Mehler) prior is the uniform ``no-information'' measure,
the Fock vacuum is its ground state, and orthogonal excitations encode
information.} The bridge is a classical limit, due in spirit to
Mehler~(1866) and made quantitative by López and Temme~\cite{LopezTemme1999}:
in the infinite-dimensional limit the uniform measure on a hypersphere
becomes the Gaussian measure, and the Gegenbauer polynomials --- which are
orthogonal with respect to the uniform (hyperspherical) measure and define
the hyperspherical harmonics --- become the Hermite polynomials.

\subsection{Uniform information is the hyperspherical measure}

Consider the surface of a $k$-sphere parameterized by angles
$\varphi_0,\dots,\varphi_{k-1}$ with $x_j=\cos\varphi_j$. Its rotation-invariant
(uniform) area element factorizes as
\begin{equation}
  \mathrm{d}\sigma(k)=\prod_{j=0}^{k-1}
  \bigl(\sqrt{1-x_j^2}\bigr)^{\,j-1}\,\mathrm{d}x_j ,
  \label{eq:areaelement}
\end{equation}
where the $j=0$ factor is special: $\mathrm{d}x_0=\sin\varphi_0\,
\mathrm{d}\varphi_0$ and the azimuthal angle $\varphi_0$ ranges over
$[0,2\pi]$, while every other angle ranges over $[0,\pi]$. The weight
$(1-x^2)^{(\alpha-1)/2}$ appearing in a single coordinate of
\eqref{eq:areaelement} is exactly the Gegenbauer weight of order
$\alpha$. Because the Gegenbauer polynomials are orthogonal with respect to
this uniform measure, the uniform measure is the natural \emph{prior}: it
expresses that we have no information, every direction on the sphere being
equally likely. Any other orthogonal state --- a definite Gegenbauer (and, in
the limit, Hermite) excitation --- encodes \emph{information}, namely the
result of integrating an unknown field's wavefunction over the hyperspherical
surface against the basis function attached to a point of that surface.

\subsection{The Gaussian limit of the hyperspherical weight}

Rescale the coordinate as $x\mapsto\sqrt{2/\alpha}\,x$ and let the dimension
$\alpha\to\infty$. The single-coordinate hyperspherical weight tends to the
Gaussian weight:
\begin{equation}
  \lim_{\alpha\to\infty}\Bigl(\sqrt{1-\tfrac{2x^2}{\alpha}}\Bigr)^{\alpha-1}
  =e^{-x^2}.
  \label{eq:gaussweight}
\end{equation}
Thus the uniform measure on an infinite-dimensional sphere is the Gaussian
measure --- the precise sense in which the Mehler ground state is the
``maximum-entropy / total-uncertainty'' prior used in
Section~\ref{sec:framework}.

\subsection{Gegenbauer \texorpdfstring{$\to$}{to} Hermite and the normalization}

The orthogonal polynomials track the same limit. With the Gegenbauer
polynomials $C_n^{(\gamma)}$ defined by the generating function
$(1-2xw+w^2)^{-\gamma}=\sum_n C_n^{(\gamma)}(x)\,w^n$, López and
Temme~\cite{LopezTemme1999} give the Hermite limit (their Eq.~(1.4), with
$\gamma=\alpha/2$):
\begin{equation}
  \lim_{\alpha\to\infty}\Bigl(\tfrac{\alpha}{2}\Bigr)^{-n/2}
  C_n^{(\alpha/2)}\!\Bigl(\sqrt{\tfrac{2}{\alpha}}\,x\Bigr)
  =\frac{1}{n!}\,H_n(x),
  \label{eq:gegtoherm}
\end{equation}
where the Hermite polynomials are
$H_n(x)=n!\sum_{k=0}^{\lfloor n/2\rfloor}\frac{(-1)^k}{k!(n-2k)!}(2x)^{n-2k}$.
The normalizations converge consistently. For integers $n,\alpha\ge0$,
\begin{equation}
  N^{(\alpha)}_n
  =\Bigl(\tfrac{\alpha}{2}\Bigr)^{\frac12-n}
  \int_{-1}^{1}\bigl[C_n^{(\alpha/2)}(x)\bigr]^2
  \bigl(\sqrt{1-x^2}\bigr)^{\alpha-1}\,\mathrm{d}x
  =\Bigl(\tfrac{\alpha}{2}\Bigr)^{\frac12-n}
  \frac{\pi\,2^{1-\alpha}\,\Gamma(n+\alpha)}
       {n!\,\bigl(n+\tfrac{\alpha}{2}\bigr)\,[\Gamma(\tfrac{\alpha}{2})]^2},
  \label{eq:gegnorm}
\end{equation}
and, in the limit, this is exactly the (squared) Hermite normalization against
the Gaussian weight:
\begin{equation}
  \lim_{\alpha\to\infty}N^{(\alpha)}_n
  =\frac{\sqrt{\pi}\,2^n}{n!}
  =\int_{-\infty}^{\infty}\Bigl[\frac{H_n(x)}{n!}\Bigr]^2 e^{-x^2}\,\mathrm{d}x .
  \label{eq:hermnorm}
\end{equation}
(The last equality uses the standard orthogonality
$\int H_m H_n\,e^{-x^2}\mathrm{d}x=\sqrt{\pi}\,2^n n!\,\delta_{mn}$.)

\subsection{Reading the dictionary}

Equations~\eqref{eq:gaussweight}--\eqref{eq:hermnorm} establish a complete
dictionary, summarized in Table~\ref{tab:dictionary}. The infinite-dimensional
hypersphere with its uniform measure \emph{is} the Gaussian/Fock setting; the
hyperspherical harmonics \emph{are} the Hermite/number-state basis; the
uniform prior \emph{is} the vacuum $\ket{0}$; and the act of measuring a
component is the projection onto a definite excitation. This is exactly the
structure the \texttt{unfer} kernel manipulates natively in
\texttt{nested\textunderscore fock\textunderscore algebra}: the vacuum
projector $\ket{0}\bra{0}$ is the uniform prior, and creating a boson in mode
$j$ records one unit of information about data channel $j$.

\begin{table}[h]
\centering
\begin{tabular}{l|l}
\textbf{Finite hypersphere ($\alpha<\infty$)} & \textbf{Infinite-dim.\ limit (Fock)} \\
\hline
uniform surface measure $\mathrm{d}\sigma$ & Gaussian measure $e^{-x^2}\mathrm{d}x$ \\
Gegenbauer weight $(1-x^2)^{(\alpha-1)/2}$ & Gaussian weight $e^{-x^2}$
  \;[Eq.~\eqref{eq:gaussweight}] \\
Gegenbauer polynomials $C_n^{(\alpha/2)}$ & Hermite polynomials $H_n/n!$
  \;[Eq.~\eqref{eq:gegtoherm}] \\
hyperspherical harmonics & number states $\ket{n}$ \\
uniform ``no-information'' prior & vacuum projector $\ket{0}\bra{0}$ \\
a definite point / harmonic & an orthogonal excitation (information) \\
\end{tabular}
\caption{The hypersphere-to-Fock dictionary underlying the Mehler prior.}
\label{tab:dictionary}
\end{table}

\section{Krylov--Hashimoto Dimensional Reduction}
\label{sec:krylov}

The generator $\bar H$ of~\eqref{eq:Hbar} (and its diagonal
surrogate~\eqref{eq:Hdiag}) exactly encodes the flow, but its nominal dimension
scales with $M$. We compress it with the SIRK machinery of
Section~\ref{sec:flowprimer}.

\subsection{The inversion-free shortcut}

The standard polynomial Krylov subspace is
\begin{equation}
  \Kry_m(\bar H,v_0)=\mathrm{span}\{v_0,\bar H v_0,\bar H^2 v_0,\dots,
  \bar H^{m-1}v_0\}.
\end{equation}
Naively, the rational construction~\eqref{eq:ratkrylov} requires solving a
linear system $(\gamma I-\bar H)y=v$ at each step. Our shortcut avoids it: by
pre-conditioning the seed as $v=(\gamma I-\bar H)^{m-1}v_0$, applying the
resolvent merely \emph{lowers the polynomial degree}, recovering the standard
Krylov subspace with no inversion. Concretely, \texttt{fock\textunderscore
sirk} realizes this by generating the sequence $w_k=(\bar H-z_kI)w_{k-1}$ with
uniform shifts $z_k=\gamma$; this is mathematically valid because $\bar H$ is
\emph{bounded} (the Mehler projector is rank one and the number operators are
diagonal). Each application of $\bar H^k$ costs only $\mathcal{O}(M)$ precisely
because the orthogonal Fock potential has \emph{no cross-terms}.

\subsection{Spectral low-pass filtering}

The Krylov projection acts as an exact spectral low-pass filter: it discards
the overfitted, localized ``spikes'' (high-frequency packet detail) and
retains the $m$ dominant macroscopic transport eigenvalues. The reduced
generator $\bar H_{\mathrm{reduced}}$ is a dense $m\times m$ matrix (e.g.\
$m=15$--$20$), and the SIRK bound~\eqref{eq:sirkbound} guarantees its error
decays like $e^{-hm}$. Generation is then a single instantaneous
$\mathcal{O}(m^2)$ matrix exponential,
\begin{equation}
  \Psi_{t=1}=e^{-i\bar H_{\mathrm{reduced}}}\,\Psi_0 .
\end{equation}
The orthonormalization of the (possibly rank-deficient) Krylov basis is handled
by \texttt{unfer}'s rank-revealing Gram whitening, which keeps the reduction
numerically stable even when several packets nearly coincide.

\section{Implementation in the \texttt{unfer} Kernel}
\label{sec:impl}

The analytical QFM module integrates into the existing \texttt{unfer} monorepo
with no new numerical primitives --- it reuses the CAS, the solver, and the
module runtime.

\subsection{Symbolic algebra (\texttt{nested\textunderscore fock\textunderscore algebra})}

The symbolic engine gains one operator, \texttt{Operator::ProjectVacuum}, the
rank-1 vacuum projector $\ket{0}\bra{0}$ implementing the Mehler prior $H_0$.
On application it keeps only the strict-vacuum component and annihilates any
state carrying a mode; it is self-adjoint and idempotent, so it slots into the
existing adjoint and Hermiticity machinery. The diagonal QFM
generator~\eqref{eq:Hdiag} is instantiated directly by the builtin
\texttt{qfm\textunderscore hamiltonian(alphas)}, which assembles
$\ket{0}\bra{0}+\sum_j\bar\alpha_j\,\hat c_j^\dagger\hat c_j$ as one projector
term plus one number operator per data point. Building the Hamiltonian
directly (rather than through \texttt{Expression::expand()}) sidesteps the
combinatorial explosion of the CAS, so $M=10^6$ data points yield $10^6$ terms
with no blow-up.

\subsection{GPU-accelerated solver (\texttt{fock\textunderscore sirk})}

The solver natively generates the Krylov sequence $w_k=(\bar H-z_kI)w_{k-1}$.
Supplying a uniform shift array realizes the inversion-free reduction of
Section~\ref{sec:krylov}; the $m\times m$ Gram matrix $\braket{w_j|w_k}$ is
assembled on the GPU via \texttt{candle-core} CUDA tensors. Because
$\hat V=\sum_j\bar\alpha_j\hat n_j$ has no cross-terms, basis generation is
embarrassingly parallel and fast, and the rank-revealing Gram whitening keeps
the highly compressed flow representation stable.

\subsection{Module execution runtime}

The whole flow-matching loop is exposed through the \texttt{unfer\textunderscore
ffi} C ABI and driven by a new Austral module (\texttt{qfm\textunderscore
module}). Through the manifest-authorized \texttt{uk\textunderscore evolve} and
\texttt{uk\textunderscore event\textunderscore probability} symbols, the module
triggers the analytical parameter assembly, the Krylov reduction, and the
single-step $\mathcal{O}(m^2)$ inference natively in-process, with the kernel
handle held in a linear type that the Austral compiler forces it to free.

\subsection{Scope: the diagonal surrogate}

The shipped builtin implements the diagonal generator~\eqref{eq:Hdiag}. Because
$\bar H_{\text{impl}}$ is diagonal, the vacuum and the data channels are
eigenstates, so $e^{-i\bar H_{\text{impl}}t}$ contributes only phases and the
Born populations are stationary --- the module's integration tests assert
exactly this conservation, not a spurious ``spread.'' The off-diagonal
differential operators $\hat h_j$ of the full generator~\eqref{eq:Hbar}, which
genuinely transport amplitude between the vacuum prior and the data channels,
are the natural next extension and fit the same SIRK pipeline unchanged.

\section{Conclusion}

We have given a self-contained derivation of an analytical Quantum Flow
Matching module for the \texttt{unfer} kernel. By encoding data as orthogonal
Fock states we turn the flow-matching normal equations into a diagonal,
$\mathcal{O}(M)$ closed-form fit, eliminating the variational PQC training of
the original QFM and its barren-plateau pathology. By compressing the
resulting bounded generator with the Shift-Invert Rational Krylov method we
obtain exponentially convergent, $\mathcal{O}(m^2)$ single-step inference. The
measure-theoretic core --- the Mehler/López--Temme limit identifying the
uniform hyperspherical measure with the Gaussian measure and Gegenbauer with
Hermite polynomials --- explains \emph{why} the Fock vacuum is the correct
``no-information'' prior and why orthogonal excitations carry information. The
result aligns with the \texttt{unfer} philosophy: rigorous probability
tracking through mathematically exact, geometry-preserving operations on
separable Hilbert spaces.

\vspace{1em}
\begin{thebibliography}{9}

\bibitem{Cui2025}
Cui, Z., Zhang, P., \& Tang, Y. (2025). \textit{Quantum Flow Matching}.
arXiv:2508.12413 [quant-ph].

\bibitem{Hashimoto2019}
Hashimoto, Y., \& Nodera, T. (2019). \textit{Shift-invert Rational Krylov
method for an operator $\phi$-function of an unbounded linear operator}.
Japan Journal of Industrial and Applied Mathematics, 36, 1--13.

\bibitem{LopezTemme1999}
López, J. L., \& Temme, N. M. (1999). \textit{Approximations of orthogonal
polynomials in terms of Hermite polynomials}. Methods and Applications of
Analysis, 6(2), 131--146.

\bibitem{Mehler1866}
Mehler, F. G. (1866). \textit{Ueber die Entwicklung einer Function von
beliebig vielen Variablen nach Laplaceschen Functionen höherer Ordnung}.
Journal für die reine und angewandte Mathematik, 66, 161--176.

\bibitem{Kopperman2026}
\textit{Beyond Real Numbers: A Hilbert-Space Exchange Format for the UniFormal
(UPL) Framework}. (2026). UniFormal architecture tutorial
(\texttt{Kopperman\textunderscore Tutorial}); Mehler uniform-measure /
Gaussian-limit foundation.

\bibitem{unfer2026}
\texttt{unfer} Kernel Architecture Documentation.
\textit{IMPLEMENTATION\textunderscore PLAN.md, ARCHITECTURE.md, PROTOCOL.md}.
(2026).

\end{thebibliography}

\end{document}
